% HOW TO: an appendix is written exactly like a chapter. \ThesisAppendices in
% main.tex counts the appendix files listed there and numbers them to match:
% on its own this prints an unlettered "APPENDIX", and next to a second
% appendix file it prints "APPENDIX A".
% Its subheadings are deliberately left out of the table of contents.
\chapter{Supplementary Material}

Supporting material that would interrupt the argument belongs here rather than
in the body: extended derivations, full instrument settings, code listings,
questionnaire text, or the signed AI Use Agreement if one is required. An
appendix is read by the few people who need it, so completeness matters more
than concision.

Each appendix is a separate file with its own \verb|\chapter| command and its
own \verb|\input| line in \verb|main.tex|, exactly like a chapter. The
\verb|\ThesisAppendices| line above those inputs is all the setup there is:
two or more appendix files are lettered A, B, C, with their tables and figures
numbered A.1, B.1 and so on, while a single appendix file prints an unlettered
``APPENDIX'' and numbers its tables and figures 1, 2, 3.

\section{An Appendix Subheading}

Appendix subheadings are styled like the ones in a chapter but are left out of
the table of contents, which lists appendix titles only. Tables and figures
inside an appendix carry the appendix letter, where there is one, and are
listed in the List of Tables and the List of Figures either way, so a reader
scanning those lists can find them without knowing they are in an appendix.

\begin{table}
  \caption{Instrument settings used throughout the study.}
  \label{tab:appendix-settings}
  \begin{tabular}{lr}
    \toprule
    Setting & Value \\
    \midrule
    Sampling rate (Hz) & 1000 \\
    Gain               &   10 \\
    \bottomrule
  \end{tabular}
\end{table}
